\section{Conclusion}\label{concl}
This paper tests a novel approach to detecting anti-competitive behavior in public procurement, applying it to auction data from São Paulo, Brazil. Inspired by \cite{kawai2023using}, we use a regression discontinuity design around narrowly decided sealed-bid auctions to identify significant incumbency advantages and strategic bidding behavior. Specifically, narrowly winning firms are more likely to have won prior auctions and to submit the final bid, indicating favoritism or collusion inconsistent with purely competitive bidding.

These findings highlight vulnerabilities in the Brazilian procurement systems, suggesting that both corruption and firm collusion may distort competition. Finally, our analysis provides a practical tool for procurement oversight bodies to flag potential anti-competitive practices and improve accountability in public contracting.
